% Pursuing Stacks (cotes 134-2 to 134-8): archive cote 134-4, site page 8 (PDF page 9), typescript page 96.
% typescript. Draft compared: ps3.tex:368-433.
% AI-transcribed by Claude Opus 5.5 (claude-opus-5-5, adaptive thinking, effort high), Message Batches API, batch ps1 (msgbatch_01Rv5hCaNTtUt4JHx1yqoHPL), 2026-11-12; not yet checked by eye.
\scanpage{134-4/8}
\mspage{96}

The basic idea of how to handle the more general situation,  namely
how to compare the categories $A_{/i^{*}(C)} = A_{/C}$ and $C$\hand{,} and show (under suitab%
le assumptions)  that their is a canonical isomorphism between their ima\xout{g}-%
ges in the localized category $W_{Cat}^{-1}(Cat) = (Hot)$,  was around since about
the moment I worked out the ``key result''.  It can be expressed by a dia-%
gram of ``maps'' in (Cat)

(1) \diagram{\begin{tikzcd} A_{/C} \ar[r] & A_{//C} & \\ & A\mathrm{x}C \ar[u] \ar[r] & C\ , \end{tikzcd}}

where $A_{//C}$ is the fibered category over $A$,  associated to the functor
\[ A^{opp} \hand{\longrightarrow} (Cat) \quad , \quad a \hand{\longmapsto} \underline{Hom}(i(a),C) \quad \hand{.} \]
Here one should be careful with the distinction  between the \uline{set}
\[ Hom(i(a),C) = Ob\ \underline{Hom}(i(a),C)\hand{,} \]
and the \uline{category} $\underline{Hom}(i(a),C)$,  both depending bi-functorially on $a$ in $A$
and $C$ in (Cat).  The former (as a presheaf on $A$ for fixed $C$)  gives rise
to $A_{/C}$ ,  a fibered category over $A$ with \uline{discrete} fibers,  whereas the
latter gives rise to $A_{//C}$ ,  which is fibered over $A$ with fibers which
need not be discrete.  Identifiying a set with the discrete category it
defines,  we get a canonical functor
\[ Hom(i(a),C) \hand{\longrightarrow} \underline{Hom}(i(a),C) \ , \]
which is very far from being an equivalence nor even a weak equivalence;
being functorial for varying $a$,  it gives rise to the first map in (1).
The second is deduced from the canonical functor
\[ C \hand{\rightarrow} \underline{Hom}(i(a),C) \quad , \]
identifying $C$ with the full subcategory of \uline{constant} functors from $i(a)$ to
$C$ .  This map is functorial in $a$,  and gives rise \inserted{again} to a cartesian functor
between the corresponding fibered categories over $A$,  the first one
(which corresponds to the constant functor $A^{opp}\hand{\rightarrow}(Cat)$ with value $C$) is
just $A\mathrm{x}C$ fibered over $A$ by $pr_{1}$ ,  hence the second\xout{\unsure{a}} arrow in (1).  The
third arrow is just $pr_{2}$ .

It seems that,  with the introduction of $A_{//C}$,  this is the first
time since the beginning of these reflections that  we are making use of
the \uline{bicategory structure} of (Cat),  namely of the notion of a morphism,
or map or ``homotopy'' (the tie with actual homotopies will be made clear
below),  between two ``maps'' namely (here) functors  $C' \hand{\rightrightarrows} C$ .  There is
of course a feeling that such a notion of homotopy should make a sense in
\xout{rather a}\inserted{a} more or less arbitrary modelizer $M$,  and that the approach dis-%
played by the diagram (1) may well generalize to more general situations

\margin{8}
